% Anexo de referencia rápida sobre archivos .bib, complementario a la
% sección "Citas y bibliografía" del capítulo "Ejemplos prácticos".
\addtocontents{toc}{ \vspace{5mm} \hrule \par}


  \chapter{Tipos de referencias}
  \label{sec:bib}
    \begin{table}[H]
	\begin{center}
	\begin{tabular}{l p{10cm} }
	 \hline
   article& Un artículo de un periódico o revista\\
  
    book& Un libro con una editorial que se indica en forma explícita. Los campos requeridos en este caso son author (autor), editor, title (título), publisher (editorial) y year (año)\\
  
    booklet& Una obra que está impresa y encuadernada, pero sin una editorial o institución patrocinadora\\
  
    conference& Lo mismo que inproceedings, incluido para compatibilidad con el lenguaje de marcación Scribe\\
  
    inbook& Una parte de un libro, que puede ser un capítulo (o sección) o un rango de páginas\\
  
    incollection& Una parte de un libro que tiene su propio título\\
  
    inproceedings& Un artículo en las actas de sesiones (proceedings) de una conferencia\\
  
    manual& Documentación técnica\\
  
    mastersthesis& Una tesis de maestría (Master thesis) o proyecto fin de carrera\\
  
    misc& Para uso cuando los demás tipos no corresponden\\
  
    phdthesis& Una tesis de doctorado (Ph D thesis)\\
  
    proceedings& Las actas de sesiones de una conferencia\\
  
    techreport& Un reporte publicado por una escuela u otra institución, usualmente numerado dentro de una serie\\
  
    unpublished& Un documento que tiene un autor y título, pero que no fue formalmente publicado\\
          \hline
	\end{tabular}
	\end{center}
	\caption{Tipos posibles de elementos en bibligrafía con natbib}
\end{table}




